 %% ---------- Vorspann
\documentclass[11pt,a4paper,twoside]{report} %article + section für hausarbeit,  report + chapter für "buch-like"
\raggedbottom
 %% ---------- Typographie
\usepackage[utf8]{inputenc}
 %% ---------- Packages und Zitierstil
\usepackage{lmodern}
\usepackage[ngerman]{babel}
 % bibliography
 % \usepackage[style=numeric,sorting=none]{biblatex} %,citestyle=authoryear
\usepackage[style=ieee,sorting=none]{biblatex}
\addbibresource{lit.bib}
\ExecuteBibliographyOptions{dashed=false}
\usepackage{subfiles}
\usepackage{tikz}
\usetikzlibrary{arrows.meta}
\usepackage{setspace}
\usepackage{graphicx}
\usepackage{scrlayer-scrpage}
\pagestyle{scrheadings}
\clearpairofpagestyles
\usepackage[a4paper,top=3cm,bottom=4cm,left=3cm,right=3cm,head=22pt]{geometry}
\usepackage{minted} % you should invoke latex with -shell-escape flag
\usepackage{listings}
%\renewcommand{\lstlistingname}{Quellcode}
%\renewcommand{\lstlistlistingname}{Quellcode}
\usepackage{csquotes}
\usepackage[md]{titlesec}
\usepackage{color}
\usepackage{xcolor}
\usepackage{amsmath}
\usepackage{mathtools}
\usepackage{fancyhdr}
\usepackage{tocbibind}
 % https://tex.stackexchange.com/questions/70610/how-to-customize-the-page-header-and-the-page-number-of-classicthesis
 % anhang mittels minted
 % \subsection{fill_database.py} 
 % \inputminted[linenos,tabsize=2,breaklines]{python}{Programmierung/fill_database.py}
\usepackage[font=small,labelfont=it]{caption}
\usepackage[hidelinks]{hyperref}
% --- Font Style
\usepackage[cmintegrals,cmbraces]{newtxmath}
\usepackage{ebgaramond-maths}
\usepackage[T1]{fontenc}
 % \usepackage[figure]{hypcap}
 % \pagestyle{fancy}
 % ---------- Header und Footer 
 % ohead should contain rightmark, which should be the current section - todo!! look up fancyhdr and renewcommand
 % for acronyms see: https://youtu.be/331YxgOJUGw?t=1541
\ohead{\rightmark}
\ofoot{\pagemark}
\def\arraystretch{1.5}%
%\discretionary{Teil1}{Teil2}{Ungetrennt}


\definecolor{mytextgrey}{RGB}{33,33,33}
\definecolor{dkgreen}{rgb}{0,0.6,0}
\definecolor{gray}{rgb}{0.5,0.5,0.5}
\definecolor{mauve}{rgb}{0.58,0,0.82}

\setlength{\footheight}{22pt}

\title{Projektarbeit}
\author{Florian Schönsee}
\date{\today}
 %% ---------- Titelseite
\begin{document}
\color{mytextgrey}
\begin{titlepage}
     % \centering   
     \hfill
     \includegraphics[width=0.4\textwidth]{bilder/beide-logos.png}
     \par\vspace{7cm}
     \raggedright
     \vspace{.5cm}
         {\scshape\Large Projektdokumentation \par}
         \vspace{1cm}
     {\LARGE \bfseries Erstellen einer Website um die drei Emotionen\\“Gelassenheit, Nostalgie, Wut” anzusprechen. \par}   
     \vspace{2cm}
     {\Large \textsc{Maria Pekar, 92 88 70 \\Scott Albers, 29 62 96\\Florian Schönsee, 94 04 85}\par}
     \vspace{0.5cm}
     \today
     \vfill
         {\small Betreuung \par}
         \textsc{Prof. Dr. Friedhelm Mündemann}
 \end{titlepage}
 

 % Vorspann, Deckblatt
\onehalfspacing
\pagenumbering{Roman}
\tableofcontents
%\newpage\null\thispagestyle{empty}\newpage
%\pagebreak
%\setcounter{page}{1}
%\addcontentsline{toc}{section}{Abkürzungsverzeichnis}
% \subfile{00Abkuerzungsverzeichnis}
 % empty page

%\newpage\null\thispagestyle{empty}\newpage
\pagebreak
\pagenumbering{Roman}
\addcontentsline{toc}{section}{Verfassererklärung}
\subfile{00_Verfassererklaerung}



\pagebreak

\pagenumbering{arabic}
%\setcounter{page}{5}
\subfile{01_Einleitung}
\pagebreak
\subfile{02_Potenzial}
\pagebreak
\subfile{03_Grundlagen}
\pagebreak
\subfile{04_Verwendung}
\pagebreak
\subfile{05_Beziehung_Konflikt}
\pagebreak
\subfile{06_Zusammenfassung}
\pagebreak
\subfile{99_Projektabschnitte}
\pagenumbering{Roman}
\setcounter{page}{2}
%\nocite{*}
\addcontentsline{toc}{chapter}{Verzeichnisse}
\addcontentsline{toc}{section}{\listfigurename}
%\listoffigures
\pagebreak
\addcontentsline{toc}{section}{\listtablename}
%\listoftables
\pagebreak
%\addcontentsline{toc}{section}{Quellcodeverzeichnis}
%\renewcommand\listoflistingscaption{Quellcodeverzeichnis}
%\listoflistings % Now typeset the list
\pagebreak
\addcontentsline{toc}{section}{Literaturverzeichnis}
\printbibliography
\pagebreak
\end{document}
